% Options for packages loaded elsewhere
\PassOptionsToPackage{unicode}{hyperref}
\PassOptionsToPackage{hyphens}{url}
%
\documentclass[
]{article}
\usepackage{amsmath,amssymb}
\usepackage{iftex}
\ifPDFTeX
  \usepackage[T1]{fontenc}
  \usepackage[utf8]{inputenc}
  \usepackage{textcomp} % provide euro and other symbols
\else % if luatex or xetex
  \usepackage{unicode-math} % this also loads fontspec
  \defaultfontfeatures{Scale=MatchLowercase}
  \defaultfontfeatures[\rmfamily]{Ligatures=TeX,Scale=1}
\fi
\usepackage{lmodern}
\ifPDFTeX\else
  % xetex/luatex font selection
\fi
% Use upquote if available, for straight quotes in verbatim environments
\IfFileExists{upquote.sty}{\usepackage{upquote}}{}
\IfFileExists{microtype.sty}{% use microtype if available
  \usepackage[]{microtype}
  \UseMicrotypeSet[protrusion]{basicmath} % disable protrusion for tt fonts
}{}
\makeatletter
\@ifundefined{KOMAClassName}{% if non-KOMA class
  \IfFileExists{parskip.sty}{%
    \usepackage{parskip}
  }{% else
    \setlength{\parindent}{0pt}
    \setlength{\parskip}{6pt plus 2pt minus 1pt}}
}{% if KOMA class
  \KOMAoptions{parskip=half}}
\makeatother
\usepackage{xcolor}
\usepackage[margin=1in]{geometry}
\usepackage{color}
\usepackage{fancyvrb}
\newcommand{\VerbBar}{|}
\newcommand{\VERB}{\Verb[commandchars=\\\{\}]}
\DefineVerbatimEnvironment{Highlighting}{Verbatim}{commandchars=\\\{\}}
% Add ',fontsize=\small' for more characters per line
\usepackage{framed}
\definecolor{shadecolor}{RGB}{248,248,248}
\newenvironment{Shaded}{\begin{snugshade}}{\end{snugshade}}
\newcommand{\AlertTok}[1]{\textcolor[rgb]{0.94,0.16,0.16}{#1}}
\newcommand{\AnnotationTok}[1]{\textcolor[rgb]{0.56,0.35,0.01}{\textbf{\textit{#1}}}}
\newcommand{\AttributeTok}[1]{\textcolor[rgb]{0.13,0.29,0.53}{#1}}
\newcommand{\BaseNTok}[1]{\textcolor[rgb]{0.00,0.00,0.81}{#1}}
\newcommand{\BuiltInTok}[1]{#1}
\newcommand{\CharTok}[1]{\textcolor[rgb]{0.31,0.60,0.02}{#1}}
\newcommand{\CommentTok}[1]{\textcolor[rgb]{0.56,0.35,0.01}{\textit{#1}}}
\newcommand{\CommentVarTok}[1]{\textcolor[rgb]{0.56,0.35,0.01}{\textbf{\textit{#1}}}}
\newcommand{\ConstantTok}[1]{\textcolor[rgb]{0.56,0.35,0.01}{#1}}
\newcommand{\ControlFlowTok}[1]{\textcolor[rgb]{0.13,0.29,0.53}{\textbf{#1}}}
\newcommand{\DataTypeTok}[1]{\textcolor[rgb]{0.13,0.29,0.53}{#1}}
\newcommand{\DecValTok}[1]{\textcolor[rgb]{0.00,0.00,0.81}{#1}}
\newcommand{\DocumentationTok}[1]{\textcolor[rgb]{0.56,0.35,0.01}{\textbf{\textit{#1}}}}
\newcommand{\ErrorTok}[1]{\textcolor[rgb]{0.64,0.00,0.00}{\textbf{#1}}}
\newcommand{\ExtensionTok}[1]{#1}
\newcommand{\FloatTok}[1]{\textcolor[rgb]{0.00,0.00,0.81}{#1}}
\newcommand{\FunctionTok}[1]{\textcolor[rgb]{0.13,0.29,0.53}{\textbf{#1}}}
\newcommand{\ImportTok}[1]{#1}
\newcommand{\InformationTok}[1]{\textcolor[rgb]{0.56,0.35,0.01}{\textbf{\textit{#1}}}}
\newcommand{\KeywordTok}[1]{\textcolor[rgb]{0.13,0.29,0.53}{\textbf{#1}}}
\newcommand{\NormalTok}[1]{#1}
\newcommand{\OperatorTok}[1]{\textcolor[rgb]{0.81,0.36,0.00}{\textbf{#1}}}
\newcommand{\OtherTok}[1]{\textcolor[rgb]{0.56,0.35,0.01}{#1}}
\newcommand{\PreprocessorTok}[1]{\textcolor[rgb]{0.56,0.35,0.01}{\textit{#1}}}
\newcommand{\RegionMarkerTok}[1]{#1}
\newcommand{\SpecialCharTok}[1]{\textcolor[rgb]{0.81,0.36,0.00}{\textbf{#1}}}
\newcommand{\SpecialStringTok}[1]{\textcolor[rgb]{0.31,0.60,0.02}{#1}}
\newcommand{\StringTok}[1]{\textcolor[rgb]{0.31,0.60,0.02}{#1}}
\newcommand{\VariableTok}[1]{\textcolor[rgb]{0.00,0.00,0.00}{#1}}
\newcommand{\VerbatimStringTok}[1]{\textcolor[rgb]{0.31,0.60,0.02}{#1}}
\newcommand{\WarningTok}[1]{\textcolor[rgb]{0.56,0.35,0.01}{\textbf{\textit{#1}}}}
\usepackage{graphicx}
\makeatletter
\def\maxwidth{\ifdim\Gin@nat@width>\linewidth\linewidth\else\Gin@nat@width\fi}
\def\maxheight{\ifdim\Gin@nat@height>\textheight\textheight\else\Gin@nat@height\fi}
\makeatother
% Scale images if necessary, so that they will not overflow the page
% margins by default, and it is still possible to overwrite the defaults
% using explicit options in \includegraphics[width, height, ...]{}
\setkeys{Gin}{width=\maxwidth,height=\maxheight,keepaspectratio}
% Set default figure placement to htbp
\makeatletter
\def\fps@figure{htbp}
\makeatother
\setlength{\emergencystretch}{3em} % prevent overfull lines
\providecommand{\tightlist}{%
  \setlength{\itemsep}{0pt}\setlength{\parskip}{0pt}}
\setcounter{secnumdepth}{-\maxdimen} % remove section numbering
% Custom definitions
% To use this customization file, insert the line "\input{custom}" in the header of the tex file.

% Formatting




% Packages

 \usepackage{amssymb,latexsym}
\usepackage{amssymb,amsfonts,amsmath,latexsym,amsthm, bm}
%\usepackage[usenames,dvipsnames]{color}
%\usepackage[]{graphicx}
%\usepackage[space]{grffile}
\usepackage{mathrsfs}   % fancy math font
% \usepackage[font=small,skip=0pt]{caption}
%\usepackage[skip=0pt]{caption}
%\usepackage{subcaption}
%\usepackage{verbatim}
%\usepackage{url}
%\usepackage{bm}
\usepackage{dsfont}
\usepackage{multirow}
%\usepackage{extarrows}
%\usepackage{multirow}
%% \usepackage{wrapfig}
%% \usepackage{epstopdf}
%\usepackage{rotating}
%\usepackage{tikz}
%\usetikzlibrary{fit}					% fitting shapes to coordinates
%\usetikzlibrary{backgrounds}	% drawing the background after the foreground


% \usepackage[dvipdfm,colorlinks,citecolor=blue,linkcolor=blue,urlcolor=blue]{hyperref}
%\usepackage[colorlinks,citecolor=blue,linkcolor=blue,urlcolor=blue]{hyperref}
%%\usepackage{hyperref}
%\usepackage[authoryear,round]{natbib}


%  Theorems, etc.

%\theoremstyle{plain}
%\newtheorem{theorem}{Theorem}[section]
%\newtheorem{corollary}[theorem]{Corollary}
%\newtheorem{lemma}[theorem]{Lemma}
%\newtheorem{proposition}[theorem]{Proposition}
%\newtheorem{condition}[theorem]{Condition}
% \newtheorem{conditions}[theorem]{Conditions}

%\theoremstyle{definition}
%\newtheorem{definition}[theorem]{Definition}
%% \newtheorem*{unnumbered-definition}{Definition}
%\newtheorem{example}[theorem]{Example}
%\theoremstyle{remark}
%\newtheorem*{remark}{Remark}
%\numberwithin{equation}{section}




% Document-specific shortcuts
\newcommand{\btheta}{{\bm\theta}}
\newcommand{\bbtheta}{{\pmb{\bm\theta}}}

\newcommand{\commentary}[1]{\ifx\showcommentary\undefined\else \emph{#1}\fi}

\newcommand{\term}[1]{\textit{\textbf{#1}}}

% Math shortcuts

% Probability distributions
\DeclareMathOperator*{\Exp}{Exp}
\DeclareMathOperator*{\TExp}{TExp}
\DeclareMathOperator*{\Bernoulli}{Bernoulli}
\DeclareMathOperator*{\Beta}{Beta}
\DeclareMathOperator*{\Ga}{Gamma}
\DeclareMathOperator*{\TGamma}{TGamma}
\DeclareMathOperator*{\Poisson}{Poisson}
\DeclareMathOperator*{\Binomial}{Binomial}
\DeclareMathOperator*{\NormalGamma}{NormalGamma}
\DeclareMathOperator*{\InvGamma}{InvGamma}
\DeclareMathOperator*{\Cauchy}{Cauchy}
\DeclareMathOperator*{\Uniform}{Uniform}
\DeclareMathOperator*{\Gumbel}{Gumbel}
\DeclareMathOperator*{\Pareto}{Pareto}
\DeclareMathOperator*{\Mono}{Mono}
\DeclareMathOperator*{\Geometric}{Geometric}
\DeclareMathOperator*{\Wishart}{Wishart}

% Math operators
\DeclareMathOperator*{\argmin}{arg\,min}
\DeclareMathOperator*{\argmax}{arg\,max}
\DeclareMathOperator*{\Cov}{Cov}
\DeclareMathOperator*{\diag}{diag}
\DeclareMathOperator*{\median}{median}
\DeclareMathOperator*{\Vol}{Vol}

% Math characters
\newcommand{\R}{\mathbb{R}}
\newcommand{\Z}{\mathbb{Z}}
\newcommand{\E}{\mathbb{E}}
\renewcommand{\Pr}{\mathbb{P}}
\newcommand{\I}{\mathds{1}}
\newcommand{\V}{\mathbb{V}}

\newcommand{\A}{\mathcal{A}}
%\newcommand{\C}{\mathcal{C}}
\newcommand{\D}{\mathcal{D}}
\newcommand{\Hcal}{\mathcal{H}}
\newcommand{\M}{\mathcal{M}}
\newcommand{\N}{\mathcal{N}}
\newcommand{\X}{\mathcal{X}}
\newcommand{\Zcal}{\mathcal{Z}}
\renewcommand{\P}{\mathcal{P}}
\newcommand{\distfn}{\mathrm{dist}} % distance measure

\newcommand{\T}{\mathtt{T}}
\renewcommand{\emptyset}{\varnothing}


% Miscellaneous commands
\newcommand{\iid}{\stackrel{\mathrm{iid}}{\sim}}
\newcommand{\matrixsmall}[1]{\bigl(\begin{smallmatrix}#1\end{smallmatrix} \bigr)}

\newcommand{\items}[1]{\begin{itemize} #1 \end{itemize}}

\newcommand{\todo}[1]{\emph{\textcolor{red}{(#1)}}}

\newcommand{\branch}[4]{
\left\{
	\begin{array}{ll}
		#1  & \mbox{if } #2 \\
		#3 & \mbox{if } #4
	\end{array}
\right.
}

% approximately proportional to
\def\app#1#2{%
  \mathrel{%
    \setbox0=\hbox{$#1\sim$}%
    \setbox2=\hbox{%
      \rlap{\hbox{$#1\propto$}}%
      \lower1.3\ht0\box0%
    }%
    \raise0.25\ht2\box2%
  }%
}
\def\approxprop{\mathpalette\app\relax}

% \newcommand{\approptoinn}[2]{\mathrel{\vcenter{
  % \offinterlineskip\halign{\hfil$##$\cr
    % #1\propto\cr\noalign{\kern2pt}#1\sim\cr\noalign{\kern-2pt}}}}}

% \newcommand{\approxpropto}{\mathpalette\approptoinn\relax}





\ifLuaTeX
  \usepackage{selnolig}  % disable illegal ligatures
\fi
\IfFileExists{bookmark.sty}{\usepackage{bookmark}}{\usepackage{hyperref}}
\IfFileExists{xurl.sty}{\usepackage{xurl}}{} % add URL line breaks if available
\urlstyle{same}
\hypersetup{
  pdftitle={Evaluation Metrics for Blocking/Entity Resolution},
  pdfauthor={STA 325: Homework 2},
  hidelinks,
  pdfcreator={LaTeX via pandoc}}

\title{Evaluation Metrics for Blocking/Entity Resolution}
\author{STA 325: Homework 2}
\date{Noah Obuya \textbar{} 1234887 \textbar{} 09/10/2024}

\begin{document}
\maketitle

\textbf{\emph{General instructions for homeworks}}: Your code must be
completely reproducible and must compile. No late homeworks will be
accepted.

\textbf{\emph{Reading}} Read the paper Binette and Steorts (2022) to get
an overview of entity resolution. You'll want to refer to this during
the course of the semester as it's meant to be a quick reference
regarding the concepts that we will be covering. For more details, refer
to the book by Christen (2012).

\textbf{\emph{Advice}}: Start early on the homeworks and it is advised
that you not wait until the day of as these homeworks are meant to be
longer and treated as case studies.

\textbf{\emph{Commenting code}} Code should be commented. See the Google
style guide for questions regarding commenting or how to write code
\url{https://google.github.io/styleguide/Rguide.xml}.

\textbf{\emph{R Markdown Test}}

\begin{enumerate}
\def\labelenumi{\arabic{enumi}.}
\setcounter{enumi}{-1}
\tightlist
\item
  Open a new R Markdown file; set the output to HTML mode and ``Knit''.
  This should produce a web page with the knitting procedure executing
  your code blocks. You can edit this new file to produce your homework
  submission.
\end{enumerate}

\textbf{Total points on assignment: 2 (reproducibility) + 23 points for the assignment = 25 total points.}

\begin{enumerate}
\def\labelenumi{\arabic{enumi}.}
\tightlist
\item
  (4 points) What are the four main challenges of entity resolution?
\item
  (4 points, 1 point each) Suppose there are 10 records in a data set.
  a.) What are the total number of brute-force comparisons needed to
  make all-to-all record comparisons? b.) Repeat this for 100 records,
  1000 records, 10,000 records. c.)
  \textcolor{blue}{As the number of records increases, comment on the growth rate of the number all-to-all record comparisons? Hint: plot all-to-all record comparisons versus total number of records}
\end{enumerate}

What do you observe about the number of comparisons that need to be
made?

\begin{enumerate}
\def\labelenumi{\arabic{enumi}.}
\setcounter{enumi}{2}
\tightlist
\item
  (9 points) Consider the following record linkage data set with
  1,000,000 total records that are matched between two databases. Assume
  that 500,000 are true matches. Assume a classified (or method) finds
  600,000 record pairs as matches, and of these 400,000 correspond as
  true matches. The number of TP + FP + TN + FN = 50,000,000.
\end{enumerate}

\begin{enumerate}
\def\labelenumi{\alph{enumi}.}
\tightlist
\item
  (4 points) Given the information above, find the following information
  in the confusion matrix: TP, FP, TN, and FN.
\item
  (1 point) Calculate the accuracy. Comment on the reliability of this
  metric for this problem.
\item
  (1 point) Calculate the precision.
\item
  (1 point) Calculate the recall.
\item
  (1 point) Calculate the f-measure.
\item
  (1 point) Comment on the reliability of the precision, recall, and
  f-measure for this problem.
\end{enumerate}

\begin{enumerate}
\def\labelenumi{\arabic{enumi}.}
\setcounter{enumi}{3}
\tightlist
\item
  (6 points) We will revisit the Italian Survey on Household and Wealth
  (SHIW) from class, which is a sample survey 383 households conducted
  by the Bank of Italy every two years (2008 and 2010). The data set is
  anonymized to remove first and last name (and other sensitive
  information).
\end{enumerate}

\begin{enumerate}
\def\labelenumi{\alph{enumi}.}
\tightlist
\item
  (0 points) Please load the data set in the way that we did in class
  and block based upon gender.
\item
  (1 point) Plot the size of the blocks and comment on how many their
  are and their relative size.
\item
  (1 point) Calculate the reduction ratio and interpret its meaning.
\item
  (2 points) Calculate the precision and recall. Interpret the meaning
  of each.
\item
  (1 point) Would this be a reasonable approach for blocking. Explain.
\item
  (1 point) Would blocking on gender be recommended for entity
  resolution. Explain.
\end{enumerate}

\newpage

\begin{Shaded}
\begin{Highlighting}[]
\FunctionTok{library}\NormalTok{(italy)}
\FunctionTok{library}\NormalTok{(assert)}
\FunctionTok{data}\NormalTok{(italy08)}
\FunctionTok{data}\NormalTok{(italy10)}
\NormalTok{knitr}\SpecialCharTok{::}\NormalTok{opts\_chunk}\SpecialCharTok{$}\FunctionTok{set}\NormalTok{(}\AttributeTok{echo =} \ConstantTok{TRUE}\NormalTok{, }
                      \AttributeTok{fig.width=}\DecValTok{4}\NormalTok{, }
                      \AttributeTok{fig.height=}\DecValTok{3}\NormalTok{, }
                      \AttributeTok{fig.align=}\StringTok{"center"}\NormalTok{)}
\FunctionTok{head}\NormalTok{(italy08)}
\end{Highlighting}
\end{Shaded}

\begin{verbatim}
##        id PARENT SEX ANASC NASCREG CIT ACOM4C STUDIO Q QUAL SETT IREG
## 1 1040021      1   2  1948      16   1      0      5 1    2    3   16
## 2 1040022     10   2  1952      16   1      0      7 1    2    3   16
## 3 1110521      1   1  1972      20   1      2      5 1    1    4   20
## 4 1110522      3   1  1935      20   1      2      2 3    6    5   20
## 5 1110523      3   2  1941      20   1      2      3 3    6    5   20
## 6  119401      1   1  1941       7   1      0      4 3    6    5    7
\end{verbatim}

\begin{Shaded}
\begin{Highlighting}[]
\FunctionTok{head}\NormalTok{(italy10)}
\end{Highlighting}
\end{Shaded}

\begin{verbatim}
##        id PARENT SEX ANASC NASCREG CIT ACOM4C STUDIO Q QUAL SETT IREG
## 1 1040021      1   2  1948      16   1      0      5 3    6    5   16
## 2 1040022     11   2  1952      16   1      0      7 1    2    3   16
## 3 1110521      1   2  1941      20   1      2      3 3    6    5   20
## 4 1110522      2   1  1935      20   1      2      2 3    6    5   20
## 5 1110523      6   1  1972      20   1      2      5 1    1    4   20
## 6  119721      1   2  1948      16   1      2      2 2    5    4   17
\end{verbatim}

\begin{Shaded}
\begin{Highlighting}[]
\NormalTok{id08 }\OtherTok{\textless{}{-}}\NormalTok{ italy08}\SpecialCharTok{$}\NormalTok{id}
\NormalTok{id10 }\OtherTok{\textless{}{-}}\NormalTok{ italy10}\SpecialCharTok{$}\NormalTok{id}
\NormalTok{id }\OtherTok{\textless{}{-}} \FunctionTok{c}\NormalTok{(italy08}\SpecialCharTok{$}\NormalTok{id, italy10}\SpecialCharTok{$}\NormalTok{id) }\CommentTok{\# combine the id}
\NormalTok{italy08 }\OtherTok{\textless{}{-}}\NormalTok{ italy08[}\SpecialCharTok{{-}}\FunctionTok{c}\NormalTok{(}\DecValTok{1}\NormalTok{)] }\CommentTok{\# remove the id}
\NormalTok{italy10 }\OtherTok{\textless{}{-}}\NormalTok{ italy10[}\SpecialCharTok{{-}}\FunctionTok{c}\NormalTok{(}\DecValTok{1}\NormalTok{)] }\CommentTok{\# remove the id }
\NormalTok{italy }\OtherTok{\textless{}{-}} \FunctionTok{rbind}\NormalTok{(italy08, italy10)}
\FunctionTok{head}\NormalTok{(italy)}
\end{Highlighting}
\end{Shaded}

\begin{verbatim}
##   PARENT SEX ANASC NASCREG CIT ACOM4C STUDIO Q QUAL SETT IREG
## 1      1   2  1948      16   1      0      5 1    2    3   16
## 2     10   2  1952      16   1      0      7 1    2    3   16
## 3      1   1  1972      20   1      2      5 1    1    4   20
## 4      3   1  1935      20   1      2      2 3    6    5   20
## 5      3   2  1941      20   1      2      3 3    6    5   20
## 6      1   1  1941       7   1      0      4 3    6    5    7
\end{verbatim}

\end{document}
